
% ════════════════════════════════════════════
% 具体分析（5.4.1.具体分析.tex，三级子块）—— 写作提示词见本文件末尾【本章节写作提示词】。
% ════════════════════════════════════════════
\subsubsection{具体分析}

针对问题四，本问属于扰动下的动态重调度与方案比较类问题，采用状态快照、枚举重调度与场景集评估相结合的方法进行求解。

首先以问题二所得调度方案为基线在 $t=5$ s 时刻建立各车辆位置与状态快照，接着刻画新增告警、随机车辆故障与通信中断及组合扰动，然后对未开始观测的告警与剩余车辆进行枚举重调度并与保持原方案策略比较，最后在 60 s 评价窗口内按完成率与最晚观测开始时间评估各场景。

具体分析流程如图\ref{fig:analysis4_flow}所示。

\begin{figure}[ht]
    \begin{center}
        \begin{tikzpicture}[x=1cm, y=0.7cm]
            \node[box] (n11) at (0, 0) {读取基线调度};
            \node[box] (n12) at (5, 0) {状态快照};
            \node[box] (n13) at (10, 0) {扰动刻画};
            \node[box] (n22) at (5, -3) {剩余告警车辆};
            \node[box] (n21) at (0, -3) {枚举重调度};
            \node[box] (n23) at (10, -3) {保持原方案};
            \node[box] (n31) at (0, -6) {场景设计};
            \node[box] (n32) at (5, -6) {指标评估};
            \node[box] (n33) at (10, -6) {收益分析};
            \draw[arrow] (n11) -- (n12);
            \draw[arrow] (n12) -- (n13);
            \draw[arrow] (n23) -- (n22);
            \draw[arrow] (n22) -- (n21);
            \draw[arrow] (n13) -- ++(0,-1.0) -| (n23);
            \draw[arrow] (n21) -- (n31);
            \draw[arrow] (n31) -- (n32);
            \draw[arrow] (n32) -- (n33);
        \end{tikzpicture}
        \caption{问题四分析流程图}
        \label{fig:analysis4_flow}
    \end{center}
\end{figure}
